% formal/hk2017-qp-precision.tex — HK2017 QP precision bounds
% Finite-precision perturbation chain and Q error bounds.
%
% Labels defined here:
%   rem:conditioning-precondition, rem:perturbation-chain,
%   lem:link-assembly, lem:link-svd, lem:link-beta0,
%   lem:link-gradient, lem:link-reduced-hessian, lem:link-eigenvalues,
%   lem:link-beta, rem:eigendirection-error,
%   lem:q-error-first-order, lem:q-correction-second-order,
%   rem:runtime-q-bound, lem:pseudoinverse-orthogonality,
%   lem:kkt-batched-defect-enclosure,
%   lem:kkt-certified-curvature-direction,
%   lem:kkt-cyclic-obstruction-inheritance,
%   cor:taylor-structure, cor:exact-correction
%
% Labels referenced but defined elsewhere:
%   rem:qp-setup, prop:q-on-affine, prop:critical-points,
%   prop:boundary-vs-interior (formal/hk2017-qp-core.tex)
%
% input by formal/main.tex — do not add \documentclass or \begin{document}
%
% ══════════════════════════════════════════════════════════════════════
% Error analysis
% ══════════════════════════════════════════════════════════════════════

\begin{remark}[Status of the finite-precision claims]\label{rem:qp-precision-status}
This file contains a mixture of proved exact-arithmetic identities,
first-order perturbation heuristics, and incomplete runtime proposals.  In
particular, the computable \(\beta\)-quantity in
\eqref{eq:eta-computable} has unproved constants and an invalid denominator
substitution, and the legacy runtime quantity in
Remark~\ref{rem:runtime-q-bound} is not an error bound for the exact
binary64-input KKT problem.  Neither may be cited as a numerical certificate.
The current proved direct certificate is instead the outward inverse-defect
argument in Lemmas~\ref{lem:kkt-verified-inverse-defect} and
\ref{lem:kkt-batched-defect-enclosure}.  This note must still be reorganized
before thesis-facing use.
\end{remark}

\begin{lemma}[Verified inverse-defect solution enclosure]%
\label{lem:kkt-verified-inverse-defect}
Let \(K\in\mathbb R^{n\times n}\), \(b\in\mathbb R^n\), let
\(\widehat x\in\mathbb R^n\) be an approximate solution, and let
\(R\in\mathbb R^{n\times n}\) be an approximate inverse.  Suppose outward
arithmetic establishes
\[
 \overline r\geq \lVert K\widehat x-b\rVert_\infty,\qquad
 \overline e\geq \lVert I-KR\rVert_\infty,\qquad
 \overline e<1.
\]
Then \(K\) is invertible and its exact solution \(x^*=K^{-1}b\) satisfies
\[
 \lVert K^{-1}\rVert_\infty
 \leq \frac{\lVert R\rVert_\infty}{1-\overline e},
 \qquad
 \lVert\widehat x-x^*\rVert_\infty
 \leq
 \rho:=
 \frac{\lVert R\rVert_\infty\,\overline r}{1-\overline e}.
\]
\end{lemma}

\begin{proof}
The matrix \(KR=I-(I-KR)\) is invertible by the Banach lemma.  Hence both
\(K\) and \(R\) are invertible, and
\[
 K^{-1}=R(KR)^{-1},\qquad
 \lVert(KR)^{-1}\rVert_\infty
 \leq (1-\overline e)^{-1}.
\]
The inverse bound follows.  Finally,
\(\widehat x-x^*=K^{-1}(K\widehat x-b)\), which gives the stated error
enclosure.
\end{proof}

\begin{lemma}[Certified positive-curvature tangent direction]%
\label{lem:kkt-certified-curvature-direction}
Let \(C\in\mathbb R^{5\times m}\) have full row rank, let
\(H\in\mathbb R^{m\times m}\) be symmetric, and let \(v\in\mathbb R^m\).
Suppose \(R\) is a right inverse of \(C\), and available outward bounds give
\[
 \|R\|_\infty\leq M,\qquad
 \|Cv\|_\infty\leq r,\qquad
 \|H\|_\infty\leq L.
\]
Put \(\eta=Mr\).  If an outward lower bound for \(v^\top Hv\) is strictly
greater than
\[
 mL\bigl(2\|v\|_\infty\eta+\eta^2\bigr),
\]
then there is a vector \(z\in\ker C\) with \(z^\top Hz>0\).
\end{lemma}

\begin{proof}
Set \(z=v-RCv\).  Then \(Cz=0\) and
\(\|z-v\|_\infty\leq\eta\).  Writing \(e=z-v\), the induced infinity norm
and \(\|w\|_1\leq m\|w\|_\infty\) give
\[
 |2v^\top He+e^\top He|
 \leq mL\bigl(2\|v\|_\infty\eta+\eta^2\bigr).
\]
The strict lower-bound hypothesis therefore implies \(z^\top Hz>0\).
\end{proof}

\begin{lemma}[Cyclic inheritance of a curvature obstruction]%
\label{lem:kkt-cyclic-obstruction-inheritance}
Let \(\tau\) and \(\sigma\) be cyclic words with distinct facet labels, and
suppose \(\tau\) injects into \(\sigma\) preserving cyclic order.  Write
\(C_\tau,H_\tau\) and \(C_\sigma,H_\sigma\) for their constraint and
quadratic matrices in the symmetric KKT convention.  If some
\(z\in\ker C_\tau\) satisfies \(z^\top H_\tau z>0\), then its zero extension
\(\bar z\) to the positions of \(\sigma\) satisfies
\[
 C_\sigma\bar z=0,\qquad
 \bar z^\top H_\sigma\bar z>0.
\]
Consequently no strictly positive feasible point on the support \(\sigma\)
can be a local maximum of the fixed-word quadratic program.
\end{lemma}

\begin{proof}
Deleting the labels outside the image of \(\tau\) leaves a cyclic rotation of
\(\tau\).  Zero extension immediately preserves all five homogeneous tangent
constraints, so \(C_\sigma\bar z=0\).

It remains only to check that choosing a different linear cut in the same
cyclic word does not change the quadratic form on \(\ker C_\tau\).  Moving
the first label \(a_0\), with coefficient \(z_0\), to the end changes the
half-quadratic pair sum \(\tfrac12z^\top H_\tau z\) by
\[
 -2z_0\sum_{j>0}z_j\,\omega_0(a_0,a_j)
 =-2z_0\,\omega_0\!\left(a_0,\sum_{j>0}z_ja_j\right)=0,
\]
because the first four rows of \(C_\tau z=0\) say
\(\sum_jz_ja_j=0\), and the omitted \(j=0\) term has
\(\omega_0(a_0,a_0)=0\).  Repeating this argument handles any cyclic
rotation.  Pair terms involving labels outside \(\tau\) vanish after zero
extension, proving equality of the two quadratic values.

Finally, \(\bar z\) is a feasible tangent direction.  At a strictly positive
feasible point, both sufficiently small signs of this direction remain
positive.  A local maximum requires the Hessian to be nonpositive on every
such tangent direction, contradicting
\(\bar z^\top H_\sigma\bar z>0\).
\end{proof}

\begin{lemma}[Batched enclosure of the exact KKT residual and inverse defect]%
\label{lem:kkt-batched-defect-enclosure}
Let \(K\) be the exact KKT matrix associated to binary64 input data, let
\(\widehat K\), \(R\), and \(\widehat x\) be binary64 matrices and vectors,
and suppose an outward binary64 matrix \(D\) satisfies
\[
 |K-\widehat K|\leq D
\]
entrywise.  Suppose also that every binary64 dot product, optionally followed
by the subtraction of a scalar \(c\), used below satisfies
\[
 \left|\operatorname{fl}\!\left(\operatorname{fl}(a^\top b)-c\right)
       -(a^\top b-c)\right|
 \leq \gamma\bigl(|a|^\top|b|+|c|\bigr)+\delta.          \tag{*}
\]
for outward binary64 constants \(\gamma<1\) and \(\delta\geq0\).

For a nonnegative dot product, let
\[
 \operatorname{updot}(a,b)
 :=
 \operatorname{up}\!\left(
   \frac{\operatorname{fl}(a^\top b)+\delta}{1-\gamma}
 \right).
\]
Apply this entrywise to matrix products, and write
\[
 \begin{aligned}
 S_R&=\operatorname{updot}(|\widehat K|,|R|),&
 T_R&=\operatorname{updot}(D,|R|),\\
 S_x&=\operatorname{updot}(|\widehat K|,|\widehat x|),&
 T_x&=\operatorname{updot}(D,|\widehat x|).
 \end{aligned}
\]
Then outward evaluation of
\[
 \begin{aligned}
 E&=
   |I-\operatorname{fl}(\widehat K R)|
   +\gamma(S_R+I)+\delta\mathbf 1+T_R,\\
 r&=
   |\operatorname{fl}(\widehat K\widehat x)-b|
   +\gamma(S_x+|b|)+\delta\mathbf 1+T_x
 \end{aligned}
\]
gives the entrywise enclosures
\[
 |I-KR|\leq E,\qquad |K\widehat x-b|\leq r.
\]
Consequently,
\(\overline e=\|E\|_\infty\) and
\(\overline r=\|r\|_\infty\) satisfy the hypotheses of
Lemma~\ref{lem:kkt-verified-inverse-defect} whenever
\(\overline e<1\).
\end{lemma}

\begin{proof}
For a nonnegative exact dot product \(s=a^\top b\), condition~\((*)\)
implies
\(\operatorname{fl}(a^\top b)\geq(1-\gamma)s-\delta\).
Thus \(s\leq\operatorname{updot}(a,b)\).
Applying~\((*)\) to the signed product and its final subtraction, and using
the corresponding nonnegative upper product, gives
\[
 \begin{aligned}
 |I-\operatorname{fl}(\widehat K R)-(I-\widehat K R)|
 &\leq\gamma(S_R+I)+\delta\mathbf 1,\\
 |\operatorname{fl}(\widehat K\widehat x)-b
          -(\widehat K\widehat x-b)|
 &\leq\gamma(S_x+|b|)+\delta\mathbf 1.
 \end{aligned}
\]
Moreover,
\[
 |(K-\widehat K)R|\leq D|R|\leq T_R.
\]
The triangle inequality now proves the defect enclosure.  Replacing \(R\)
by \(\widehat x\) gives the residual enclosure.  Outward row sums give the
claimed infinity-norm bounds.
\end{proof}

\begin{remark}[Binary64 implementation contract for the batched enclosure]%
\label{rem:kkt-batched-binary64-contract}
The development implementation uses
\(\varepsilon=2^{-52}\), \(s=2n\varepsilon\),
\(\gamma=\operatorname{up}(s/(1-s))\), and
\(\delta=\operatorname{up}(4n\,2^{-1074})\).
Thus it budgets approximately \(4n\) unit roundoffs for an augmented dot
product that uses at most \(n\) multiplications, \(n-1\) additions, and the
final subtraction from \(b\) or \(I\).  The dynamic
matrix path of the pinned \texttt{nalgebra 0.33.3} dependency dispatches to
\texttt{matrixmultiply 0.3.10}; for the current KKT sizes its ordinary and
fused multiply-add kernels evaluate one inner-dimension block.  The separate
\texttt{nalgebra 0.35.0} dependency supplies only the Bunch--Kaufman
factorization.  Fused
multiply-add uses fewer roundings than the separately rounded model.
Multiplication by the GEMM scaling factor \(1\) is exact.

The runtime checks gradual underflow once before each route with opaque runtime
operands, covering both flush-to-zero outputs and denormals-are-zero inputs.
The floating-point environment is required to remain unchanged during that
single-threaded route.  If gradual underflow is unavailable, the route bypasses
all floating-point certificates---including short-word and curvature
rejections---and exact-resolves the original candidate stream.  The batched
predicate itself returns \textsc{Indeterminate}
if any required value is non-finite, or if the resulting defect norm is not
strictly below one.  A dependency, target, compiler-arithmetic, maximum-word-
length, or matrix-multiplication change requires this implementation contract
to be checked again.  The exact theorem is
Lemma~\ref{lem:kkt-batched-defect-enclosure}; the dependency inspection and
exact-binary64 experiments are implementation evidence, not part of its proof.
\end{remark}

\begin{corollary}[Certified beta, \(Q\), and action signs from one KKT enclosure]%
\label{cor:kkt-beta-q-from-xi}
Use the symmetric KKT convention
\[
 \begin{pmatrix}H&C^\top\\ C&0\end{pmatrix}
 \binom{\beta}{\lambda}
 =
 \binom{0}{d},
 \qquad
 d=(0,0,0,0,1)^\top,
\]
and let \(\xi\) be the last component of \(\lambda\).  Under the hypotheses of
Lemma~\ref{lem:kkt-verified-inverse-defect}, every component of the exact
\(\beta^*\) and \(\xi^*\) lies within \(\rho\) of its computed value.  Thus
\[
 \widehat\beta_i-\rho>0\quad\hbox{for all \(i\)}
\]
certifies strict positivity, while
\(\widehat\beta_i+\rho\leq0\) for some \(i\) certifies failure of strict
positivity.

Moreover,
\[
 2Q(\beta^*)+\xi^*=0,
\]
so outward endpoint rounding gives
\[
 Q(\beta^*)\in
 \left[-\frac{\widehat\xi+\rho}{2},
       -\frac{\widehat\xi-\rho}{2}\right].
\]
For \(0<q_{\rm L}\leq Q(\beta^*)\leq q_{\rm U}\), the action lies in
\[
 \frac1{2q_{\rm U}}
 \leq \mathcal A
 \leq \frac1{2q_{\rm L}}.
\]
\end{corollary}

\begin{proof}
The component enclosures follow from the infinity-norm bound.  Multiplying
the stationarity equation \(H\beta^*+C^\top\lambda^*=0\) by
\(\beta^{*\top}\), and using \(C\beta^*=d\), gives
\(\beta^{*\top}H\beta^*+\lambda^{*\top}d=2Q(\beta^*)+\xi^*=0\).
The remaining statements follow by monotonicity.
\end{proof}
%
% Structure:
%   (A) Perturbation chain: a_i -> C -> V -> H' -> gamma_j -> beta
%   (B) Q error bounds (existing, from earlier work)
%   (C) beta > 0 certification (the main open problem)


% ── A. Perturbation chain ────────────────────────────────────────────

\begin{remark}[Conditioning precondition]\label{rem:conditioning-precondition}
All bounds in the perturbation chain
(Lemmas~\ref{lem:link-assembly}--\ref{lem:link-beta})
have $\sigma_{\min}(C)$ in their denominators.
When $\sigma_{\min}(C)$ is small:
\begin{itemize}
  \item The null-space basis error $\|\delta V\|$ grows as
    $O(\varepsilon_{\mathrm{mach}} / \sigma_{\min}(C))$
    (Lemma~\ref{lem:link-svd});
  \item The particular solution error $\|\delta\beta_0\|$ grows as
    $O(\varepsilon_{\mathrm{mach}} / \sigma_{\min}(C)^2)$
    (Lemma~\ref{lem:link-beta0});
  \item The eigenvalue perturbation $\varepsilon_\gamma$ grows as
    $O(\varepsilon_{\mathrm{mach}} / \sigma_{\min}(C))$,
    potentially exceeding $|\gamma_j|$ for all eigenvalues.
\end{itemize}
The f64 solver checks $\tilde\sigma_{\min}(C) > \varepsilon_C$
as a precondition.
If the check fails, the solver returns \textsc{Indeterminate}
for the entire $\sigma$-node, because no downstream bound
can certify anything.

The threshold $\varepsilon_C$ is not a single derived value.
Rather, $\sigma_{\min}(C)$ enters all downstream bounds, and the
solver computes those bounds using the observed $\tilde\sigma_{\min}(C)$.
If all resulting bounds are too large to certify
(eigenvalue signs, $\beta > 0$, or Q~error), the outcome is
\textsc{Indeterminate} regardless of any fixed threshold.
The precondition $\tilde\sigma_{\min}(C) > \varepsilon_C$ is
a fast short-circuit: skip the chain when the bounds are
guaranteed to be useless.

% [GAP - The value of eps_C is determined empirically: choose
% the smallest value such that no certified (TRUE/FALSE) verdict
% is overturned by exact arithmetic.  The perturbation chain
% bounds are valid for any sigma_min(C) > 0; the question is
% whether they are tight enough to certify.]
\end{remark}


\begin{remark}[Perturbation chain overview]\label{rem:perturbation-chain}
The f64 solver computes a chain of quantities, each depending
on the previous.
At each link, errors from f64 arithmetic and from perturbed
inputs accumulate.
We bound the error at each link using standard backward
stability results for the f64 primitives (SVD, eigendecomposition)
and perturbation theory (Weyl, Davis--Kahan).

The chain, with the key bound at each link:
\[
  a_i
  \;\xrightarrow{\text{assemble}}\;
  C, H
  \;\xrightarrow{\text{SVD}}\;
  V, \beta_0
  \;\xrightarrow{V^\top H V}\;
  H'
  \;\xrightarrow{\text{eig}}\;
  \gamma_j
  \;\xrightarrow{(H')^{-1} g}\;
  \alpha^*
  \;\xrightarrow{\beta_0 + V\alpha^*}\;
  \beta^*.
\]
\end{remark}


\begin{lemma}[Link 1: Assembly]\label{lem:link-assembly}
$H$ and $C$ are assembled from the dual vertices $a_i$ by
$O(m^2)$ elementary f64 operations
($\omega_0$ evaluations, copies).
The computed $\tilde C$ and $\tilde H$ satisfy
\[
  \|\tilde C - C\| \leq c_1\, \varepsilon_{\mathrm{mach}}\, \|C\|,
  \qquad
  \|\tilde H - H\| \leq c_2\, \varepsilon_{\mathrm{mach}}\, \|H\|,
\]
where $c_1, c_2$ are small constants depending on~$m$
(at most $O(m)$), and $\varepsilon_{\mathrm{mach}} \approx 1.1 \times 10^{-16}$.
% [TODO: JÖRN - verify constants; the assembly involves
% a subtraction (omega_0 = q1*p2 - q2*p1) which could cause
% cancellation if a_i are nearly aligned.  For our polytopes,
% a_i are O(1) with no near-cancellation, so this is safe.]

In our setting, $\|C\|$ and $\|H\|$ are $O(1)$
(dual vertices $a_i$ are $O(1)$),
so $\|\delta C\|$ and $\|\delta H\|$ are $O(\varepsilon_{\mathrm{mach}})$.
\end{lemma}


\begin{lemma}[Link 2: Null-space basis via SVD]\label{lem:link-svd}
% Citation: Golub \& Van Loan (2013), §8.6 (Computing the SVD) for backward stability.
%   Exact theorem number within §8.6 not confirmed — verify at library if needed.
The SVD of~$\tilde C$ is backward stable: the computed
$\tilde U \tilde\Sigma \tilde V_{\mathrm{full}}^\top$
is the exact SVD of $\tilde C + \delta C_{\mathrm{SVD}}$
where $\|\delta C_{\mathrm{SVD}}\| \leq c_3\, \varepsilon_{\mathrm{mach}}\, \|\tilde C\|$.

The null-space basis $\tilde V$ (the last $k$ right singular vectors)
approximates the exact null space of~$C$.
The angle between the computed and exact null spaces is bounded by
the \textbf{Davis--Kahan $\sin\theta$ theorem}:
\[
  \|\sin\Theta(\tilde V, V)\|
  \leq \frac{\|\delta C_{\mathrm{total}}\|}
       {\sigma_p(C) - \sigma_{p+1}(C)}
  = \frac{\|\delta C_{\mathrm{total}}\|}
       {\sigma_{\min}(C)},
\]
where $\delta C_{\mathrm{total}}$ combines assembly and SVD errors,
$\sigma_p(C) = \sigma_{\min}(C)$ is the smallest nonzero singular
value ($C$ has full row rank~$p = 5$),
and $\sigma_{p+1}(C) = 0$ (the gap to the null space).

% [GAP - The sin(Theta) bound requires sigma_min(C) >> 0.
% If sigma_min(C) is small, the null space of C is
% ill-conditioned: a small perturbation of C changes V
% significantly.  In the EHZ setting, sigma_min(C) depends
% on the geometry of the dual vertices a_{sigma(i)}.
% Empirically, sigma_min(C) < 0.01 is associated with
% boundary-case beta (margin ~ 1e-12), so ill-conditioned
% null spaces coincide with cases the solver should return
% INDETERMINATE anyway.]
\end{lemma}


\begin{lemma}[Link 2b: Particular solution]\label{lem:link-beta0}
The particular solution $\beta_0 = C^+ d$ (minimum-norm solution
of $C\beta = d$) is computed from the SVD of~$C$.
Write $\tilde\beta_0 = \tilde C^+ d$ for the computed value.
Since $C$ has full row rank, $C^+ = C^\top(C C^\top)^{-1}$
and $\|C^+\| = 1/\sigma_{\min}(C)$.
The standard pseudoinverse perturbation bound
for full-row-rank matrices gives:
% Citation: Wedin (1973), "Perturbation theory for pseudo-inverses", BIT 13, 217--232.
%   Also Stewart (1977), SIAM Review 19, 634--662 (survey). The sqrt(2) bound is Wedin's central result.
%   NOTE: This is NOT the Wedin (1972) BIT 12:99--111 paper in bibliography.bib (that one is
%   about singular value perturbation bounds). A separate bib entry for Wedin 1973 may be needed.
\[
  \|\delta\beta_0\|
  = \|\tilde\beta_0 - \beta_0\|
  \leq \frac{c_6\,\|\delta C_{\mathrm{total}}\| \cdot \|d\|}
       {\sigma_{\min}(C)^2}
  + O(\text{second order}),
\]
where $c_6$ is a small constant.
Using $\|\delta C_{\mathrm{total}}\| = O(\varepsilon_{\mathrm{mach}} \cdot \|C\|)$
from Lemmas~\ref{lem:link-assembly}--\ref{lem:link-svd}:
\begin{equation}\label{eq:beta0-perturbation}
  \|\delta\beta_0\|
  = O\!\left(
    \frac{\varepsilon_{\mathrm{mach}} \cdot \|C\| \cdot \|d\|}
         {\sigma_{\min}(C)^2}
  \right).
\end{equation}
The $1/\sigma_{\min}(C)^2$ scaling is one power worse than the
null-space error ($1/\sigma_{\min}(C)$ from Lemma~\ref{lem:link-svd}),
because $\beta_0$ depends on the inverses of the singular values of~$C$,
not just their separation from zero.

We also record $\|\beta_0\| \leq \|C^+\| \cdot \|d\|
= \|d\| / \sigma_{\min}(C)$, which enters subsequent bounds.
\end{lemma}


\begin{lemma}[Link 2c: Reduced gradient]\label{lem:link-gradient}
The reduced gradient $g = V^\top H \beta_0$ is computed as
$\tilde g = \tilde V^\top \tilde H\, \tilde\beta_0$.
Write $\delta g = \tilde g - g$.  Expanding to first order:
\begin{align*}
  \delta g
  &= (V + \delta V)^\top (H + \delta H)(\beta_0 + \delta\beta_0)
     - V^\top H \beta_0 \\
  &= V^\top H\, \delta\beta_0
     + \delta V^\top H\, \beta_0
     + V^\top \delta H\, \beta_0
     + O(\text{second order}).
\end{align*}
Each term is bounded by:
\begin{enumerate}
  \item $\|V^\top H\, \delta\beta_0\|
    \leq \|H\| \cdot \|\delta\beta_0\|$
    \quad (Lemma~\ref{lem:link-beta0}).
  \item $\|\delta V^\top H\, \beta_0\|
    \leq \|H\| \cdot \|\delta V\| \cdot \|\beta_0\|$
    \quad (Lemma~\ref{lem:link-svd}).
  \item $\|V^\top \delta H\, \beta_0\|
    \leq \|\delta H\| \cdot \|\beta_0\|$
    \quad (Lemma~\ref{lem:link-assembly}).
\end{enumerate}
Substituting:
$\|\delta\beta_0\| = O(\varepsilon_{\mathrm{mach}} \cdot \|C\| / \sigma_{\min}(C)^2)$,\;
$\|\delta V\| = O(\varepsilon_{\mathrm{mach}} / \sigma_{\min}(C))$,\;
$\|\beta_0\| \leq \|d\| / \sigma_{\min}(C)$,\;
$\|\delta H\| = O(\varepsilon_{\mathrm{mach}})$:
\begin{equation}\label{eq:gradient-perturbation}
  \|\delta g\|
  = O\!\left(
    \frac{(\|H\| \cdot \|C\| + \|H\| + 1) \cdot
          \varepsilon_{\mathrm{mach}} \cdot \|d\|}
         {\sigma_{\min}(C)^2}
  \right).
\end{equation}
In the EHZ setting ($\|H\|$, $\|C\|$, $\|d\|$ all $O(1)$),
this simplifies to $\|\delta g\| = O(\varepsilon_{\mathrm{mach}} / \sigma_{\min}(C)^2)$.
\end{lemma}


\begin{lemma}[Link 3: Reduced Hessian]\label{lem:link-reduced-hessian}
The reduced Hessian $H' = V^\top H V$ is computed as
$\tilde H' = \tilde V^\top \tilde H\, \tilde V$.
Write $\tilde V = V + \delta V$ where $\|\delta V\|$ is bounded
by Lemma~\ref{lem:link-svd}. Then:
\begin{align*}
  \tilde H'
  &= (V + \delta V)^\top (H + \delta H) (V + \delta V) \\
  &= V^\top H V
     + V^\top H\, \delta V
     + \delta V^\top H V
     + V^\top \delta H\, V
     + O(\|\delta V\|^2, \|\delta H\| \cdot \|\delta V\|) \\
  &= H' + \Delta H',
\end{align*}
where the first-order perturbation satisfies
\[
  \|\Delta H'\|
  \leq 2\,\|H\|\,\|\delta V\| + \|\delta H\| + O(\text{second order}).
\]
Using $\|\delta H\| = O(\varepsilon_{\mathrm{mach}})$
and $\|\delta V\| = O(\varepsilon_{\mathrm{mach}} / \sigma_{\min}(C))$
from Lemmas~\ref{lem:link-assembly}--\ref{lem:link-svd}:
\[
  \|\Delta H'\|
  = O\!\left(
    \frac{\|H\|\, \varepsilon_{\mathrm{mach}}}{\sigma_{\min}(C)}
  \right).
\]
\end{lemma}


\begin{lemma}[Link 4: Eigenvalues of the reduced Hessian]%
\label{lem:link-eigenvalues}
% Citation: Weyl (1912), Math. Annalen 71, 441--479. Modern: Horn \& Johnson (2013),
%   Matrix Analysis, Theorem 4.3.1. For singular values: follows from eigenvalue version.
By Weyl's eigenvalue perturbation theorem,
the computed eigenvalues $\tilde\gamma_j$ of~$\tilde H'$
satisfy
\[
  |\tilde\gamma_j - \gamma_j|
  \leq \|\tilde H' - H'\|
  \leq \|\Delta H'\| + c_4\, \varepsilon_{\mathrm{mach}}\, \|\tilde H'\|,
\]
where the second term accounts for f64 rounding in the
eigendecomposition itself (backward stable).
Combining with Lemma~\ref{lem:link-reduced-hessian}:
\begin{equation}\label{eq:eigenvalue-perturbation}
  |\tilde\gamma_j - \gamma_j|
  = O\!\left(
    \frac{\|H\|\, \varepsilon_{\mathrm{mach}}}{\sigma_{\min}(C)}
  \right).
\end{equation}

\textbf{Eigenvalue sign reliability:}
if $|\tilde\gamma_j| > \varepsilon_\gamma$
where
$\varepsilon_\gamma = c_5 \cdot \|H\|\, \varepsilon_{\mathrm{mach}} / \sigma_{\min}(C)$,
then $\operatorname{sign}(\tilde\gamma_j) = \operatorname{sign}(\gamma_j)$.
The trinary classification threshold is derived, not heuristic.
\end{lemma}


\begin{remark}[Unproved Link 5 proposal for $\beta > 0$]%
\label{lem:link-beta}
Assume all retained eigenvalues of~$\tilde H'$ are negative:
$\tilde\gamma_j < -\varepsilon_\gamma$ for $j = 1, \ldots, k$
(no null eigenvalues; the semidefinite case with null directions
is treated separately via the LP search in the f64 algorithm).
By Lemma~\ref{lem:link-eigenvalues}, the exact eigenvalues
satisfy $\gamma_j < 0$, so $H'$ is negative definite
and the unique critical point is
$\alpha^* = -(H')^{-1} g$ with $\beta^* = \beta_0 + V\alpha^*$.

Let $w_1, \ldots, w_k$ be orthonormal eigenvectors of~$H'$
(so $H' w_j = \gamma_j w_j$, $g_j = w_j^\top g$,
$\alpha^*_j = -g_j / \gamma_j$).
The computed $\tilde\alpha = -(\tilde H')^{-1} \tilde g$ has
first-order error:
% Citation: GVL (2013), §2.6 (Sensitivity of Square Systems) confirmed.
%   Higham §7.2 is componentwise; the normwise identity is in §7.1. Use GVL §2.6 as primary.
\begin{equation}\label{eq:alpha-perturbation}
  \delta\alpha
  = \tilde\alpha - \alpha^*
  = -(H')^{-1}\bigl(\Delta H' \cdot \alpha^* + \delta g\bigr)
  + O(\text{second order}).
\end{equation}

\textbf{Direction-dependent amplification.}
In the eigenbasis of~$H'$, the $j$-th component is:
\begin{equation}\label{eq:alpha-component-error}
  (\delta\alpha)_j
  = -\frac{e_j + f_j}{\gamma_j},
  \qquad
  e_j = w_j^\top \Delta H'\, \alpha^*,
  \quad
  f_j = w_j^\top \delta g.
\end{equation}
The amplification $1/|\gamma_j|$ makes the error large for
eigenvalues near zero.

\textbf{Full $\beta$ error.}
The error in component~$i$ of $\beta$ ($i = 1, \ldots, m$) has three sources:
\begin{equation}\label{eq:beta-error-decomposition}
  (\delta\beta)_i
  = \underbrace{\sum_j (Vw_j)_i \cdot (\delta\alpha)_j}_{%
      \text{critical-point shift}}
  + \underbrace{(\delta V \cdot \alpha^*)_i}_{%
      \text{null-space rotation}}
  + \underbrace{(\delta\beta_0)_i}_{%
      \text{particular-solution shift}}
  + O(\text{second order}).
\end{equation}

\textbf{Componentwise bound.}
Each term in~\eqref{eq:beta-error-decomposition} is bounded by:
\begin{itemize}
  \item $|e_j| \leq \|\Delta H'\| \cdot \|\alpha^*\|$
    (Cauchy--Schwarz; Lemma~\ref{lem:link-reduced-hessian}).
  \item $|f_j| \leq \|\delta g\|$
    (Lemma~\ref{lem:link-gradient}).
  \item $|(\delta V \cdot \alpha^*)_i|
    \leq \|\delta V\| \cdot \|\alpha^*\|$
    (Lemma~\ref{lem:link-svd}).
  \item $|(\delta\beta_0)_i| \leq \|\delta\beta_0\|$
    (Lemma~\ref{lem:link-beta0}).
\end{itemize}
Combining (for each component $i = 1, \ldots, m$):
\begin{equation}\label{eq:beta-certification-bound}
  |(\delta\beta)_i|
  \leq \bigl(\|\Delta H'\| \cdot \|\alpha^*\| + \|\delta g\|\bigr)
       \sum_{j=1}^{k} \frac{|(Vw_j)_i|}{|\gamma_j|}
  + \|\delta V\| \cdot \|\alpha^*\|
  + \|\delta\beta_0\|.
\end{equation}

\textbf{Computable proxy.}
Since exact quantities ($V$, $w_j$, $\gamma_j$, $\alpha^*$)
are unavailable, the f64 solver substitutes computed values.
For retained eigenvalues ($|\tilde\gamma_j| > \varepsilon_\gamma$),
Lemma~\ref{lem:link-eigenvalues} gives
$|\gamma_j| \geq |\tilde\gamma_j| - \varepsilon_\gamma > 0$.
Define the computable bound (for each $i = 1, \ldots, m$;
the sum runs over retained eigenvalues only,
i.e.\ those with $|\tilde\gamma_j| > \hat E_{\Delta H'}$):
\begin{equation}\label{eq:eta-computable}
  \eta_i
  = \bigl(\hat E_{\Delta H'} \|\tilde\alpha\|
    + \hat E_{\delta g}\bigr)
    \sum_{j=1}^{k}
      \frac{|(\tilde V \tilde w_j)_i|}
           {|\tilde\gamma_j|}
  + \hat E_{\delta V} \|\tilde\alpha\|
  + \hat E_{\delta\beta_0},
\end{equation}
where $\eta_i = \infty$ if any retained eigenvalue has
$|\tilde\gamma_j| \leq \hat E_{\Delta H'}$ (perturbation
dominates the eigenvalue; first-order analysis invalid).
Using $|\tilde\gamma_j|$ instead of the proved lower bound
$|\tilde\gamma_j|-\varepsilon_\gamma$ is not valid for an upper error bound.
The displayed implementation proxy therefore remains heuristic even before
the unresolved constants below are considered.
where the constants are computable upper bounds from
Lemmas~\ref{lem:link-svd}--\ref{lem:link-gradient}:
\begin{align*}
  \hat E_{\Delta H'}
  &= c_{\Delta H'}\,
     \frac{\|\tilde H\|\, \varepsilon_{\mathrm{mach}}}
          {\tilde\sigma_{\min}(C)},
  &
  \hat E_{\delta g}
  &= c_{\delta g}\,
     \frac{\|\tilde H\| \cdot \|\tilde C\| \cdot
           \varepsilon_{\mathrm{mach}}}
          {\tilde\sigma_{\min}(C)^2}, \\
  \hat E_{\delta V}
  &= c_{\delta V}\,
     \frac{\varepsilon_{\mathrm{mach}}}
          {\tilde\sigma_{\min}(C)},
  &
  \hat E_{\delta\beta_0}
  &= c_{\delta\beta_0}\,
     \frac{\|\tilde C\| \cdot
           \varepsilon_{\mathrm{mach}}}
          {\tilde\sigma_{\min}(C)^2}.
\end{align*}
Tildes denote computed (f64) values; the substitution
of computed for exact quantities introduces $O(\varepsilon_{\mathrm{mach}})$
relative errors, absorbed into the constants.
All quantities ($\tilde V$, $\tilde w_j$, $\tilde\gamma_j$,
$\tilde\alpha$, $\tilde H$, $\tilde C$,
$\tilde\sigma_{\min}(C)$) are available from the solver.

\medskip\noindent
\textbf{Trinary $\beta > 0$ certification.}
\begin{itemize}
  \item \textsc{True}
    ($\tilde\beta_i - \eta_i > 0$ for all~$i$):
    then $\beta^*_i \geq \tilde\beta_i - |(\delta\beta)_i|
    \geq \tilde\beta_i - \eta_i > 0$.
    Interior solution certified.
  \item \textsc{False}
    ($\tilde\beta_i + \eta_i < 0$ for some~$i$):
    then $\beta^*_i \leq \tilde\beta_i + \eta_i < 0$.
    Infeasible $\beta$ certified.
  \item \textsc{Indeterminate}: otherwise.
    Cannot determine the sign of~$\beta^*_i$.
\end{itemize}

% [GAP - The constants c_{DeltaH'}, c_{delta g}, c_{delta V},
% c_{delta beta_0} are not derived from first principles.
% In the implementation, we use c = m^2 as a safety factor.
% Empirical validation: zero violations on 364 I1 natural polytope
% problems and all synthetic well-conditioned problems.
% 39 violations on null-eigenvalue cases (|gamma_j| < eps_gamma)
% where the bound does not apply — these are INDETERMINATE.]
\end{remark}


\begin{remark}[Per-eigendirection error structure]%
\label{rem:eigendirection-error}
From~\eqref{eq:alpha-component-error},
$|(\delta\alpha)_j| \leq (|e_j| + |f_j|) / |\gamma_j|$
where $|e_j|, |f_j| = O(\varepsilon_{\mathrm{mach}})$
(Lemmas~\ref{lem:link-reduced-hessian}--\ref{lem:link-gradient}).
This gives
\begin{equation}\label{eq:eigendirection-scaling}
  |(\delta\alpha)_j|
  = O\!\left(\frac{\varepsilon_{\mathrm{mach}}}{|\gamma_j|}\right).
\end{equation}

\textbf{Empirical confirmation}
(364 natural polytope problems in case~I1, all $\gamma_j < 0$):
the product $|(\delta\alpha)_j| \cdot |\gamma_j|$ is
$10^{-17}$--$10^{-16}$ across 15 orders of magnitude of~$|\gamma_j|$
($10^{-5}$ to $10^0$), confirming~\eqref{eq:eigendirection-scaling}
with proportionality constant $\approx \varepsilon_{\mathrm{mach}}$.
The $\eta_i$ bound~\eqref{eq:eta-computable} uses $c = m^2$
as a safety factor; zero violations on all well-conditioned problems.

\textbf{Limitation:}
When $|\gamma_j| \lesssim \varepsilon_\gamma$
($= O(\|H\| \varepsilon_{\mathrm{mach}} / \sigma_{\min}(C))$,
Lemma~\ref{lem:link-eigenvalues}),
the first-order analysis breaks down:
$\|\Delta H'\| \gtrsim |\gamma_j|$, so the perturbation
is not small relative to the eigenvalue.
Empirically, $|(\delta\alpha)_j| = O(1)$ for $|\gamma_j| \approx 10^{-15}$.
These cases must be classified as \textsc{Indeterminate}.

% Empirical data: verify-numerics experiment, filter_poly_diverse,
% 364 I1 problems from 307 polytopes (max 8 facets).
% Diagnostic fields: proj_delta_alpha, proj_eigenvalues in results.jsonl.
\end{remark}


% ── B. Q error bounds (existing) ─────────────────────────────────────

% The lemmas below were proven before the perturbation chain above.
% They bound |Q(tilde beta) - Q(beta^*)| using the KKT residual r,
% without going through the chain.  This is a complementary approach:
% the chain gives componentwise understanding, while the residual
% bound gives a single scalar bound on Q error.
%
% The KKT system M x = b where M = [H, C^T; C, 0] is the augmented
% form of the per-sigma problem.  The residual r = M tilde x - b
% measures how well the computed solution satisfies the KKT conditions.

\begin{lemma}[First-order Q error bound]\label{lem:q-error-first-order}
Let $Q(\beta) = \tfrac{1}{2}\beta^\top H \beta$ where
$H \in \mathbb{R}^{m \times m}$ is symmetric.
Let $C \in \mathbb{R}^{p \times m}$ have full row rank,
$d \in \mathbb{R}^p$, and
$M = \bigl[\begin{smallmatrix} H & C^\top \\ C & 0 \end{smallmatrix}\bigr]$.
Suppose $\beta^* > 0$ solves the interior KKT system $M(\beta^*,\lambda^*)^\top = (0,d)^\top$,
and $(\tilde\beta, \tilde\lambda)$ is a computed solution with residual
$r = M(\tilde\beta,\tilde\lambda)^\top - (0,d)^\top$.
Write $\delta\beta = \tilde\beta - \beta^*$ and
$r_\lambda = C\tilde\beta - d = C\,\delta\beta$ for the
constraint residual (the last $p$ components of~$r$).
Then
\[
  |Q(\tilde\beta) - Q(\beta^*)|
  \;\leq\;
  \frac{\|H\|\,\|\beta^*\|}{\sigma_{\min}(C)}\,\|r\|
  \;+\; \tfrac{1}{2}\,\|H\|\,\|\delta\beta\|^2,
\]
where $\|\cdot\|$ denotes the spectral norm for matrices and the
Euclidean norm for vectors, and $\sigma_{\min}(C)$ is the smallest
singular value of~$C$.
\end{lemma}

\begin{proof}
Taylor expansion at $\beta^*$:
\begin{equation}\label{eq:q-taylor}
  Q(\tilde\beta) - Q(\beta^*)
  = (H\beta^*)^\top \delta\beta + \tfrac{1}{2}\,\delta\beta^\top H\,\delta\beta.
\end{equation}

\emph{First-order term.}
Since $\beta^* > 0$, the interior KKT stationarity condition gives
$H\beta^* + C^\top\lambda^* = 0$, i.e.,
$H\beta^* = -C^\top\lambda^*$. Substituting:
\begin{equation}\label{eq:first-order-rewrite}
  (H\beta^*)^\top\delta\beta
  = -(C^\top\lambda^*)^\top\delta\beta
  = -\lambda^{*\top} C\,\delta\beta
  = -\lambda^{*\top} r_\lambda.
\end{equation}
Since $C$ has full row rank~$p$, $C^\top \in \mathbb{R}^{m \times p}$
is injective (its null space is trivial). Therefore
$\|C^\top y\| \geq \sigma_{\min}(C^\top)\|y\| = \sigma_{\min}(C)\|y\|$
for all $y \in \mathbb{R}^p$. Applying this to $C^\top\lambda^* = -H\beta^*$:
\[
  \sigma_{\min}(C)\,\|\lambda^*\|
  \leq \|C^\top\lambda^*\|
  = \|H\beta^*\|
  \leq \|H\|\,\|\beta^*\|.
\]
Thus $\|\lambda^*\| \leq \|H\|\,\|\beta^*\| / \sigma_{\min}(C)$.
Since $r_\lambda$ is a sub-vector of $r$ (Euclidean norm),
$\|r_\lambda\| \leq \|r\|$.
Therefore
\[
  |(H\beta^*)^\top\delta\beta|
  = |\lambda^{*\top} r_\lambda|
  \leq \|\lambda^*\|\,\|r_\lambda\|
  \leq \frac{\|H\|\,\|\beta^*\|}{\sigma_{\min}(C)}\,\|r\|.
\]

\emph{Second-order term.}
$|\tfrac{1}{2}\,\delta\beta^\top H\,\delta\beta|
  \leq \tfrac{1}{2}\,\|H\|\,\|\delta\beta\|^2$.

Combining with \eqref{eq:q-taylor} by the triangle inequality:
\[
  |Q(\tilde\beta) - Q(\beta^*)|
  \leq \frac{\|H\|\,\|\beta^*\|}{\sigma_{\min}(C)}\,\|r\|
  + \tfrac{1}{2}\,\|H\|\,\|\delta\beta\|^2. \qedhere
\]
\end{proof}


\begin{lemma}[Q correction is second-order]\label{lem:q-correction-second-order}
Under the same assumptions as Lemma~\ref{lem:q-error-first-order},
define the \emph{Q correction} as $\Delta Q = \tilde\lambda^\top r_\lambda$,
where $\tilde\lambda$ is the computed Lagrange multiplier and
$r_\lambda = C\tilde\beta - d$.
Then the corrected Q value
$Q_{\mathrm{corr}} = Q(\tilde\beta) + \Delta Q$ satisfies
\[
  Q_{\mathrm{corr}} - Q(\beta^*)
  = \delta\lambda^\top r_\lambda
    + \tfrac{1}{2}\,\delta\beta^\top H\,\delta\beta,
\]
where $\delta\lambda = \tilde\lambda - \lambda^*$.
In particular,
$|Q_{\mathrm{corr}} - Q(\beta^*)|
  \leq \|\delta\lambda\|\,\|r_\lambda\|
    + \tfrac{1}{2}\,\|H\|\,\|\delta\beta\|^2$.
\end{lemma}

\begin{proof}
From \eqref{eq:q-taylor} and \eqref{eq:first-order-rewrite}:
$Q(\tilde\beta) - Q(\beta^*) = -\lambda^{*\top} r_\lambda
+ \tfrac{1}{2}\,\delta\beta^\top H\,\delta\beta$. Adding the correction:
\begin{align*}
  Q_{\mathrm{corr}} - Q(\beta^*)
  &= Q(\tilde\beta) + \tilde\lambda^\top r_\lambda - Q(\beta^*) \\
  &= -\lambda^{*\top} r_\lambda + \tilde\lambda^\top r_\lambda
    + \tfrac{1}{2}\,\delta\beta^\top H\,\delta\beta \\
  &= (\tilde\lambda - \lambda^*)^\top r_\lambda
    + \tfrac{1}{2}\,\delta\beta^\top H\,\delta\beta \\
  &= \delta\lambda^\top r_\lambda
    + \tfrac{1}{2}\,\delta\beta^\top H\,\delta\beta.
\end{align*}
The bound follows from the triangle inequality and
$|\delta\beta^\top H\,\delta\beta| \leq \|H\|\,\|\delta\beta\|^2$.
\end{proof}


\begin{remark}[Legacy runtime Q diagnostic]\label{rem:runtime-q-bound}
At runtime, the exact $\beta^*$ is unknown.
The first-order term of Lemma~\ref{lem:q-error-first-order}
can be computed using the solver output $\tilde\beta$
as a proxy for~$\beta^*$:
\[
  E_1 = \frac{\|H\|\,\|\tilde\beta\|}{\sigma_{\min}(C)}\,\|r\|.
\]
This does not in general bound the first-order part of the Q~error: replacing
\(\beta^*\) by \(\tilde\beta\) introduces an uncontrolled term, and an ordinary
rounded residual is not an exact residual enclosure.
The second-order term $\tfrac{1}{2}\|H\|\,\|\delta\beta\|^2$
cannot be computed directly (since $\delta\beta$ is unknown),
It may be small on well-conditioned problems, but the current implementation
does not prove the required conditioning and rounding hypotheses.

% [GAP - The claim "||delta_beta|| is small when ||r|| is small" requires
% bounding ||delta_beta|| in terms of ||r|| via the conditioning of M.
% The naive bound ||delta_beta|| <= ||M^+|| * ||r|| is valid but
% ||M^+|| can be very large (up to 10^16 in experiments). A tighter
% analysis via the block structure of M is needed to make this rigorous.]

All quantities in $E_1$ are available from the solver:
$\|H\|$ from the eigendecomposition (largest $|\lambda_i(H)|$),
$\|\tilde\beta\|$ from the solution vector,
$\|r\|$ from the KKT residual, and
$\sigma_{\min}(C)$ from the SVD of~$C$
(already computed by the constraint solver).

% [TODO: JÖRN - Empirical validation shows max ratio
%   |Q_err| / E_1 = 0.196 across 477 test problems from 15 matrix families.
%   The constant 1 in the bound is not tight; the proof gives an upper
%   bound, not an equality.  Whether a tighter constant can be proven
%   (e.g., involving dimension m) is open.]
\end{remark}


\begin{lemma}[Pseudoinverse orthogonality]\label{lem:pseudoinverse-orthogonality}
Let $M$ be symmetric and $\tilde x = M^+ b$ the pseudoinverse solution.
Let $x^*$ satisfy $M x^* = b$ with $x^* \in \mathrm{col}(M)$.
Define $\delta x = \tilde x - x^*$ and $r = M\tilde x - b$.
Then $\delta x^\top M\, \delta x = 0$.
\end{lemma}

\begin{proof}
Since $M$ is symmetric, $M^+$ maps into $\mathrm{col}(M) = \mathrm{row}(M)$,
so $\tilde x \in \mathrm{col}(M)$. By assumption $x^* \in \mathrm{col}(M)$,
hence $\delta x \in \mathrm{col}(M)$.
From $M\tilde x = b + r$ and $M x^* = b$, we get $M\,\delta x = r$.
The residual satisfies $r \perp \mathrm{col}(M)$ (standard pseudoinverse property).
Therefore $\delta x^\top r = 0$ (since $\delta x \in \mathrm{col}(M)$), and
\[
  \delta x^\top M\,\delta x = \delta x^\top r = 0. \qedhere
\]
\end{proof}


\begin{corollary}[Taylor structure of Q error]\label{cor:taylor-structure}
Under the assumptions of Lemma~\ref{lem:q-error-first-order} and
Lemma~\ref{lem:pseudoinverse-orthogonality}, the uncorrected Q~error satisfies
\[
  Q(\tilde\beta) - Q(\beta^*)
  = (H\beta^*)^\top\delta\beta + \tfrac{1}{2}\,\delta\beta^\top H\,\delta\beta
\]
with $(H\beta^*)^\top\delta\beta = -\delta\beta^\top H\,\delta\beta$,
i.e., the first-order term is exactly $-2$ times the second-order term.
Consequently, $Q(\tilde\beta) - Q(\beta^*) = -\tfrac{1}{2}\,\delta\beta^\top H\,\delta\beta$.
\end{corollary}

\begin{proof}
Expanding $\delta x^\top M\,\delta x = 0$ in blocks:
\[
  \delta\beta^\top H\,\delta\beta + 2\,\delta\beta^\top C^\top\delta\lambda = 0.
\]
From $x^{*\top} M\,\delta x = 0$ (same argument with $x^*$ in place of $\delta x$):
\[
  \beta^{*\top} H\,\delta\beta + d^\top\delta\lambda + \lambda^{*\top} r_\lambda = 0,
\]
using $C\beta^* = d$ and $C\,\delta\beta = r_\lambda$.
Subtracting gives
$(H\beta^*)^\top\delta\beta - \delta\beta^\top H\,\delta\beta
= d^\top\delta\lambda + \lambda^{*\top} r_\lambda - 2\,r_\lambda^\top\delta\lambda$.
% [GAP - The identity (Hβ*)^T δβ = -δβ^T H δβ is verified empirically
% (ratio = 2.0000 across 95 rank-deficient cases) but the algebraic
% proof above yields extra terms d^T δλ + λ*^T r_λ - 2 r_λ^T δλ
% that need to cancel. This cancellation follows from x*^T r = 0
% (since x* ∈ col(M)), giving β*^T r_β + λ*^T r_λ = 0, combined
% with the block expansion. Full derivation needs Jörn.]
\end{proof}


\begin{corollary}[Exact correction in exact arithmetic]\label{cor:exact-correction}
Under the same assumptions,
\[
  Q_{\mathrm{corr}} - Q(\beta^*)
  = \delta\lambda^\top r_\lambda + \tfrac{1}{2}\,\delta\beta^\top H\,\delta\beta.
\]
From Lemma~\ref{lem:pseudoinverse-orthogonality},
$\tfrac{1}{2}\,\delta\beta^\top H\,\delta\beta = -r_\lambda^\top\delta\lambda$
(expanding $\delta x^\top M\,\delta x = 0$).
Since $\delta\lambda^\top r_\lambda$ and $r_\lambda^\top\delta\lambda$ are
both scalars and equal, $Q_{\mathrm{corr}} - Q(\beta^*) = 0$.
\end{corollary}
